
\documentclass[test,letterpaper]{cs20}
\usepackage{enumerate}
\usepackage{tikz}
\usepackage{pgf}
\usepackage{hyperref}
\usepackage{amsmath}

\begin{document}

\header{5}{Final Exam (Make-up)}


\begin{problem}

Let $A$, $B$, and $C$ be sets.

\subproblem Prove that $(A \setminus B) \setminus C = A \setminus (B \cup C)$. (Hint: show $\subseteq$ in both directions.)

\subproblem Using part (a), prove the cardinality formula
\[
|(A \setminus B) \setminus C| = |A| - |A \cap (B \cup C)|
\]
for finite sets $A$, $B$, $C$. You may use without proof that for finite sets $X \subseteq Y$, $|Y \setminus X| = |Y| - |X|$, and the standard identity $A \setminus S = A \setminus (A \cap S)$ for any $S$.

\subproblem State a condition on $A$, $B$, $C$ that characterizes exactly when $(A \setminus B) \setminus C = A$. (No justification required.)

\begin{solution}
\subsolution ($\subseteq$) Let $x \in (A \setminus B) \setminus C$. Then $x \in A \setminus B$ and $x \notin C$, i.e., $x \in A$, $x \notin B$, and $x \notin C$. So $x \notin B \cup C$, hence $x \in A \setminus (B \cup C)$.

($\supseteq$) Let $x \in A \setminus (B \cup C)$. Then $x \in A$, $x \notin B$, and $x \notin C$. Then $x \in A \setminus B$ and $x \notin C$, so $x \in (A \setminus B) \setminus C$. $\blacksquare$

\subsolution By part (a), $(A \setminus B) \setminus C = A \setminus (B \cup C)$. Using $A \setminus S = A \setminus (A \cap S)$ with $S = B \cup C$, and then the given fact about cardinality of subtraction:
\[
|A \setminus (B \cup C)| = |A \setminus (A \cap (B \cup C))| = |A| - |A \cap (B \cup C)|. \quad \blacksquare
\]

\subsolution $(A \setminus B) \setminus C = A$ iff $A \cap (B \cup C) = \emptyset$, i.e., $A$ is disjoint from both $B$ and $C$.
\end{solution}

\pagebreak

\end{problem}


\begin{problem}

Let $A$ and $B$ be (possibly infinite) sets. For functions $f: A \to B$ and $g: B \to A$, write $g \circ f: A \to A$ for the composition $(g \circ f)(x) = g(f(x))$. Let $\mathrm{id}_A: A \to A$ denote the identity function on $A$.

Consider the claim: \emph{For all functions $f: A \to B$ and $g: B \to A$, if $g \circ f = \mathrm{id}_A$, then $f$ is a bijection and $g = f^{-1}$.}

\subproblem Disprove the claim by giving explicit finite sets $A$, $B$ and explicit functions $f$, $g$ satisfying the hypothesis but not the conclusion. Verify each property.

\subproblem Prove the \emph{strengthened claim:} For all functions $f: A \to B$ and $g: B \to A$, if $g \circ f = \mathrm{id}_A$ and $f \circ g = \mathrm{id}_B$, then $f$ is a bijection and $g = f^{-1}$.

\begin{solution}
\subsolution Take $A = \{1, 2\}$, $B = \{1, 2, 3\}$, with $f(1) = 1$, $f(2) = 2$, and $g(1) = 1$, $g(2) = 2$, $g(3) = 1$. Verify the hypothesis: $(g \circ f)(1) = g(1) = 1$ and $(g \circ f)(2) = g(2) = 2$, so $g \circ f = \mathrm{id}_A$. But $f$ is not surjective (no element of $A$ maps to $3 \in B$), so $f$ is not a bijection. Hence the conclusion fails.

\subsolution Assume $g \circ f = \mathrm{id}_A$ and $f \circ g = \mathrm{id}_B$.

\emph{$f$ is injective.} Suppose $f(x_1) = f(x_2)$ for some $x_1, x_2 \in A$. Apply $g$ to both sides: $g(f(x_1)) = g(f(x_2))$, i.e., $(g \circ f)(x_1) = (g \circ f)(x_2)$, i.e., $x_1 = x_2$.

\emph{$f$ is surjective.} For any $y \in B$, take $x = g(y) \in A$. Then $f(x) = f(g(y)) = (f \circ g)(y) = y$.

So $f$ is a bijection. Finally, $g = f^{-1}$: for every $y \in B$, $f(g(y)) = y$ (from the second hypothesis), which is exactly the defining property of the inverse. $\blacksquare$

The counterexample violates the strengthened hypothesis: $f \circ g \neq \mathrm{id}_B$, since $(f \circ g)(3) = f(g(3)) = f(1) = 1 \neq 3$.
\end{solution}

\pagebreak

\end{problem}


\begin{problem}

Define a relation $R$ on $\mathbb{Z}$ by
\[
a \mathbin{R} b \iff 3 \mid (a - b).
\]

\subproblem Prove that $R$ is an equivalence relation on $\mathbb{Z}$ (i.e., $R$ is reflexive, symmetric, and transitive).

\subproblem Describe the equivalence classes of $R$. How many distinct classes are there, and what does each one look like?

\begin{solution}
\subsolution
\emph{Reflexive.} For any $a \in \mathbb{Z}$, $a - a = 0$ and $3 \mid 0$, so $a \mathbin{R} a$.

\emph{Symmetric.} Suppose $a \mathbin{R} b$, i.e., $3 \mid (a - b)$, so $a - b = 3k$ for some $k \in \mathbb{Z}$. Then $b - a = -3k = 3(-k)$, so $3 \mid (b - a)$, i.e., $b \mathbin{R} a$.

\emph{Transitive.} Suppose $a \mathbin{R} b$ and $b \mathbin{R} c$, so $a - b = 3k$ and $b - c = 3m$ for integers $k, m$. Then $a - c = (a - b) + (b - c) = 3k + 3m = 3(k + m)$, so $3 \mid (a - c)$, i.e., $a \mathbin{R} c$. $\blacksquare$

\subsolution The equivalence class of $a$ is $\{b \in \mathbb{Z} : 3 \mid (a - b)\}$, which is the set of integers congruent to $a$ modulo $3$.

There are exactly $\boxed{3}$ distinct equivalence classes:
\begin{itemize}
\item $[0] = \{\ldots, -6, -3, 0, 3, 6, \ldots\}$ (integers divisible by $3$).
\item $[1] = \{\ldots, -5, -2, 1, 4, 7, \ldots\}$ (integers $\equiv 1 \pmod 3$).
\item $[2] = \{\ldots, -4, -1, 2, 5, 8, \ldots\}$ (integers $\equiv 2 \pmod 3$).
\end{itemize}
Every integer lies in exactly one of these three classes.
\end{solution}

\pagebreak

\end{problem}


\begin{problem}

Six people --- including Alice and Bob --- are to be seated in a row of $6$ chairs.

\subproblem In how many distinct orderings can the six people be seated, with no further restrictions?

\subproblem In how many of those orderings do Alice and Bob sit next to each other (in either order)?

\begin{solution}
\subsolution The number of arrangements of $6$ distinct people in a row is $6! = 720$.

\subsolution Treat Alice and Bob as a single "block." Then there are $5$ items to arrange (the AB-block plus the other $4$ people), giving $5! = 120$ arrangements. Within the block, Alice and Bob can be in $2$ internal orders (A--B or B--A). By the product rule, the count is $5! \cdot 2 = 240$.
\end{solution}

\pagebreak

\end{problem}


\begin{problem}

\subproblem Verify that $1^3 + 2^3 + \cdots + n^3 = \left(\dfrac{n(n+1)}{2}\right)^2$ for $n = 1, 2, 3$ by direct computation.

\subproblem Prove by induction on $n$ that $1^3 + 2^3 + \cdots + n^3 = \left(\dfrac{n(n+1)}{2}\right)^2$ for every integer $n \geq 1$.

\begin{solution}
\subsolution
\begin{itemize}
\item $n = 1$: LHS $= 1^3 = 1$ and RHS $= (1 \cdot 2 / 2)^2 = 1$. $\checkmark$
\item $n = 2$: LHS $= 1 + 8 = 9$ and RHS $= (2 \cdot 3 / 2)^2 = 9$. $\checkmark$
\item $n = 3$: LHS $= 1 + 8 + 27 = 36$ and RHS $= (3 \cdot 4 / 2)^2 = 36$. $\checkmark$
\end{itemize}

\subsolution Let $P(n) := \left(\sum_{i=1}^n i^3 = \left(\frac{n(n+1)}{2}\right)^2\right)$. Proof for all $n \geq 1$ by induction.

\emph{Base case ($n = 1$).} $1^3 = 1 = \left(\frac{1 \cdot 2}{2}\right)^2$, so $P(1)$ holds.

\emph{Inductive step.} Assume $P(k)$. Compute:
\[
\sum_{i=1}^{k+1} i^3 = \sum_{i=1}^k i^3 + (k+1)^3 = \left(\frac{k(k+1)}{2}\right)^2 + (k+1)^3.
\]
Factor $(k+1)^2$:
\[
= (k+1)^2 \left[\frac{k^2}{4} + (k+1)\right] = (k+1)^2 \cdot \frac{k^2 + 4(k+1)}{4} = (k+1)^2 \cdot \frac{k^2 + 4k + 4}{4} = (k+1)^2 \cdot \frac{(k+2)^2}{4} = \left(\frac{(k+1)(k+2)}{2}\right)^2.
\]
Hence $P(k+1)$ holds, and by induction $P(n)$ holds for all $n \geq 1$. $\blacksquare$
\end{solution}

\pagebreak

\end{problem}


\begin{problem}

Let $S$ be the smallest set of integer pairs satisfying:
\begin{itemize}
\item \textbf{Base:} $(1, 1) \in S$.
\item \textbf{C1:} If $(x, y) \in S$, then $(x + 1,\ 3 y) \in S$.
\end{itemize}

\subproblem Show that $(4, 27) \in S$ by exhibiting an explicit derivation. List each rule application and the resulting pair.

\subproblem Prove by structural induction that every $(x, y) \in S$ satisfies $y = 3^{x - 1}$.

\begin{solution}
\subsolution Derivation of $(4, 27)$:
\begin{enumerate}
\item Base: $(1, 1) \in S$.
\item C1: $(2, 3) \in S$.
\item C1: $(3, 9) \in S$.
\item C1: $(4, 27) \in S$. $\checkmark$
\end{enumerate}
Sanity check: $3^{4 - 1} = 27$. $\checkmark$

\subsolution \textbf{Base case.} $(1, 1)$: $3^{1 - 1} = 1 = y$. $\checkmark$

\textbf{Inductive step.} Assume $(x, y) \in S$ has $y = 3^{x - 1}$. Then C1 produces $(x + 1, 3y)$. Compute:
\[
3 y = 3 \cdot 3^{x - 1} = 3^x = 3^{(x + 1) - 1}.
\]
By structural induction, every $(x, y) \in S$ satisfies $y = 3^{x - 1}$. $\blacksquare$
\end{solution}

\pagebreak

\end{problem}


\begin{problem}

A factory produces parts with defect rate $\pi = 0.2$. Each part is independently checked by two quality-control tests:
\begin{itemize}
\item Test $A$: probability of flagging given defective is $0.8$, probability of flagging given non-defective is $0.1$.
\item Test $B$: probability of flagging given defective is $0.9$, probability of flagging given non-defective is $0.2$.
\end{itemize}
Conditional on the true defect status, the two test results are independent.

\subproblem Compute $P(\text{defective} \mid A^+ \cap B^+)$, the probability a part is defective given both tests flag it.

\subproblem Compute $P(\text{defective} \mid A^- \cap B^+)$, the probability a part is defective given test $A$ does not flag but test $B$ does.

\begin{solution}
Let $D$ = defective ($P(D) = 0.2$, $P(\overline{D}) = 0.8$).

\subsolution
\[
P(A^+ B^+ \mid D) = 0.8 \cdot 0.9 = 0.72, \quad P(A^+ B^+ \mid \overline{D}) = 0.1 \cdot 0.2 = 0.02.
\]
\[
P(D \mid A^+ B^+) = \frac{0.72 \cdot 0.2}{0.72 \cdot 0.2 + 0.02 \cdot 0.8} = \frac{0.144}{0.144 + 0.016} = \frac{0.144}{0.160} = \frac{9}{10}.
\]

\subsolution
\[
P(A^- B^+ \mid D) = 0.2 \cdot 0.9 = 0.18, \quad P(A^- B^+ \mid \overline{D}) = 0.9 \cdot 0.2 = 0.18.
\]
\[
P(D \mid A^- B^+) = \frac{0.18 \cdot 0.2}{0.18 \cdot 0.2 + 0.18 \cdot 0.8} = \frac{0.036}{0.036 + 0.144} = \frac{0.036}{0.180} = \frac{1}{5}.
\]
\end{solution}

\pagebreak

\end{problem}


\begin{problem}

You roll a fair eight-sided die (with faces $1, 2, \ldots, 8$) repeatedly until you see a value in $\{1, 8\}$. Let $N$ be the number of rolls needed. Since the number of rolls is a geometric random variable with success probability $1/4$, $\mathbb{E}[N] = 4$. You may use this fact without proof.

\subproblem Let $T$ be the value of the \emph{final} roll (so $T \in \{1, 8\}$). Compute $\mathbb{E}[T]$.

\subproblem Let $S$ be the sum of the values of all rolls (including the final roll that produced the value in $\{1, 8\}$). Compute $\mathbb{E}[S]$.

\begin{solution}
\subsolution By symmetry, at each roll the values $1$ and $8$ are equally likely (each with probability $1/8$, independently of past rolls). Hence the first of $\{1, 8\}$ to appear is equally likely to be either, so
\[
\mathbb{E}[T] = \tfrac{1}{2}(1 + 8) = \tfrac{9}{2} = 4.5.
\]

\subsolution Decompose $S = (\text{sum of first } N - 1 \text{ rolls}) + T$. The first $N - 1$ rolls are conditionally uniform on $\{2, 3, 4, 5, 6, 7\}$ (six values, since they were not in $\{1, 8\}$), with conditional mean
\[
\frac{2 + 3 + 4 + 5 + 6 + 7}{6} = \frac{27}{6} = 4.5.
\]
By linearity of expectation, using $\mathbb{E}[N] = 4$,
\[
\mathbb{E}[\text{sum of first } N-1] = (\mathbb{E}[N] - 1) \cdot 4.5 = 3 \cdot 4.5 = 13.5.
\]
Therefore
\[
\mathbb{E}[S] = 13.5 + 4.5 = 18.
\]
\end{solution}

\pagebreak

\end{problem}


\begin{problem}

Consider the directed graph $G$ on vertex set $V = \{1, 2, 3, 4, 5, 6, 7\}$ with edge set
\[
E(G) = \{(1, 3),\ (2, 3),\ (3, 4),\ (4, 5),\ (5, 4),\ (4, 6),\ (4, 7)\}.
\]
(So $(u, v) \in E$ means there is a directed edge from $u$ to $v$.)

\subproblem Identify the strongly connected components (SCCs) of $G$. List each SCC as a set of vertices, and describe the \emph{condensation DAG} of $G$ --- the directed acyclic graph obtained by collapsing each SCC of $G$ to a single super-node, with an edge from super-node $X$ to super-node $Y$ whenever some edge of $G$ goes from a vertex in $X$ to a vertex in $Y$ (with $X \neq Y$).

\subproblem Determine the minimum number of additional directed edges that must be added to $G$ to make it strongly connected, and exhibit one valid set of additions. Justify the minimum using the structure of the condensation DAG. (Hint: count sources and sinks of the condensation; when the condensation has at least two vertices, the answer is $\max(\#\text{sources}, \#\text{sinks})$.)

\begin{solution}
\subsolution Trace cycles in $G$: the only cycle is $4 \to 5 \to 4$, giving an SCC $\{4, 5\}$. Vertices $1$, $2$, $3$, $6$, $7$ each have no return path and are singleton SCCs.

SCCs: $\{1\},\ \{2\},\ \{3\},\ \{4, 5\},\ \{6\},\ \{7\}$ ($6$ in total).

Condensation DAG: project each edge of $G$ onto its SCC endpoints (skipping edges internal to an SCC):
\begin{itemize}
\item $(1, 3) \rightsquigarrow \{1\} \to \{3\}$.
\item $(2, 3) \rightsquigarrow \{2\} \to \{3\}$.
\item $(3, 4) \rightsquigarrow \{3\} \to \{4, 5\}$.
\item $(4, 6) \rightsquigarrow \{4, 5\} \to \{6\}$.
\item $(4, 7) \rightsquigarrow \{4, 5\} \to \{7\}$.
\item $(4, 5)$ and $(5, 4)$ are internal to $\{4, 5\}$ --- they do not appear in the condensation.
\end{itemize}
So the condensation DAG has the $5$ edges
\[
\{1\} \to \{3\}, \quad \{2\} \to \{3\}, \quad \{3\} \to \{4, 5\}, \quad \{4, 5\} \to \{6\}, \quad \{4, 5\} \to \{7\}.
\]

\subsolution In the condensation DAG:
\begin{itemize}
\item \emph{Sources} (no incoming edges): $\{1\}$ and $\{2\}$. So $\#\text{sources} = 2$.
\item \emph{Sinks} (no outgoing edges): $\{6\}$ and $\{7\}$. So $\#\text{sinks} = 2$.
\end{itemize}
By the standard result, the minimum number of edges to add is $\max(2, 2) = 2$.

\emph{One valid set of $2$ additions:} the edges $(6, 1)$ and $(7, 2)$ (in the original graph $G$). Verify the augmented graph is strongly connected:
\begin{itemize}
\item Cycle $1 \to 3 \to 4 \to 6 \to 1$ exists, so $\{1, 3, 4, 5, 6\}$ are mutually reachable.
\item Cycle $2 \to 3 \to 4 \to 7 \to 2$ exists, so $\{2, 3, 4, 5, 7\}$ are mutually reachable.
\item These two cycles share vertex $3$ (and $4, 5$), so all $7$ vertices lie in one big SCC.
\end{itemize}

\emph{Lower bound argument:} every source SCC must gain at least one incoming edge from outside (else it remains a source after augmentation), and every sink SCC must gain at least one outgoing edge. By a pairing argument (each added edge can serve at most one source and one sink), at least $\max(\#\text{sources}, \#\text{sinks})$ edges are necessary. $\blacksquare$
\end{solution}

\pagebreak

\end{problem}


\begin{problem}

Student Alex proposes the following algorithm for finding a minimum-weight spanning tree (MST) on a connected weighted graph $G$ with $n$ vertices: \emph{Pick a starting vertex $v_0$ and add it to the tree. Maintain a ``current vertex'' $u$, initially set to $v_0$. Repeat $n - 1$ times: among edges incident to $u$ whose other endpoint is not yet in the tree, add the one of minimum weight; then update $u$ to be that newly added vertex.} (Alex only ever considers edges leaving the most recently added vertex, not edges leaving other tree vertices.)

\subproblem Draw a connected weighted graph on $4$ vertices with $5$ edges, and choose a starting vertex $v_0$, on which Alex's algorithm successfully produces an MST. Specify the edge weights, the chosen $v_0$, and report the total MST weight.

\subproblem Draw a connected weighted graph on $4$ vertices, and choose a starting vertex $v_0$, on which Alex's algorithm fails --- either the algorithm cannot complete $n - 1$ steps (the current vertex runs out of usable edges) or the set of $n - 1$ edges produced is not a minimum-weight spanning tree. Specify the weights, the chosen $v_0$, and explain in $1$--$2$ sentences exactly why the algorithm fails.

\subproblem Suggest a one-line modification to Alex's algorithm that would make it correct (i.e., always produce an MST on connected weighted graphs). Which well-known MST algorithm does it become?

\begin{solution}
\subsolution Take $V = \{A, B, C, D\}$ with $v_0 = A$ and edges:
\[
AB = 1, \quad BC = 2, \quad CD = 3, \quad AD = 10, \quad AC = 20.
\]
Trace Alex's algorithm starting at $A$:
\begin{itemize}
\item $u = A$. Edges from $A$ to non-tree vertices: $AB(1), AC(20), AD(10)$. Cheapest: $AB$. Add. $u \leftarrow B$.
\item $u = B$. Edges from $B$ to non-tree vertices ($\{C, D\}$): $BC(2)$ only. Add. $u \leftarrow C$.
\item $u = C$. Edges from $C$ to non-tree vertices ($\{D\}$): $CD(3)$ only. Add. Done.
\end{itemize}
Output: $\{AB, BC, CD\}$, the path $A - B - C - D$, total weight $6$ --- the MST (any other spanning tree includes an edge of weight $\geq 10$).

\subsolution Take $V = \{A, B, C, D\}$ with $v_0 = A$ and edges:
\[
AB = 1, \quad AC = 2, \quad AD = 3, \quad BC = 100.
\]
The graph is connected (everything reaches everything else through $A$). Trace Alex's algorithm:
\begin{itemize}
\item $u = A$. Cheapest edge from $A$ to a non-tree vertex: $AB(1)$. Add. $u \leftarrow B$.
\item $u = B$. Edges from $B$ to non-tree vertices ($\{C, D\}$): $BC(100)$ only. (No edge $BD$ exists.) Add. $u \leftarrow C$.
\item $u = C$. Edges from $C$ to non-tree vertices ($\{D\}$): \emph{none} (no edge $CD$ exists). Algorithm cannot proceed.
\end{itemize}
The algorithm halts after $2$ edges and never reaches $D$ --- not a spanning tree. The issue: from $C$, no edge to $D$ exists, but $A$ (already in the tree) has the edge $AD(3)$ available --- Alex's algorithm refuses to use it because $A$ is no longer the current vertex.

\subsolution \emph{Modified algorithm:} For $n - 1$ steps, add the edge of minimum weight among all edges with exactly one endpoint anywhere in the current tree (the cheapest edge across the whole frontier, not just the last-added vertex). This is \textbf{Prim's algorithm}. Considering the whole frontier guarantees that any reachable non-tree vertex remains accessible through some tree vertex, so the tree always grows and after $n - 1$ steps spans all $n$ vertices.
\end{solution}

\pagebreak

\end{problem}

\end{document}
